\section{Introducción}

El objetivo principal de esta práctica es aplicar conceptos fundamentales de seguridad informática mediante la simulación de un ataque controlado a un servidor remoto, empleando diversas técnicas y herramientas de auditoría de redes y descifrado de contraseñas.

El flujo de trabajo se dividió en varias fases estratégicas. Inicialmente, se llevó a cabo una fase de reconocimiento en la que cada integrante analizó su red local para familiarizarse con el uso de comandos de red e identificar la topología y los dispositivos conectados. Posteriormente, el análisis se dirigió hacia un servidor objetivo (18.219.229.216), utilizando la herramienta \texttt{nmap} para evadir restricciones de firewall, identificar puertos abiertos, e inferir los servicios y el sistema operativo en ejecución.

El núcleo de la práctica consistió en obtener acceso a la cuenta \texttt{purple} del servidor. Dado que las contraseñas reales se encontraban dentro de un diccionario con millones de registros, fue indispensable aplicar técnicas de filtrado y depuración utilizando comandos de UNIX (como \texttt{grep} y \texttt{split}) guiados por una pista interceptada. Esta reducción permitió ejecutar un ataque de diccionario distribuido y eficiente mediante la herramienta \texttt{hydra}. Finalmente, tras descifrar la credencial, se accedió al servidor mediante el protocolo SSH para dejar constancia de la incursión, cumpliendo así con los lineamientos establecidos sin comprometer la integridad de la infraestructura ni la de otros equipos. Ademas de esto, se respondieron preguntas de teoría relacionadas con la práctica, reforzando el aprendizaje de los conceptos aplicados.